\chapter{Editorial and Writing Guidelines}

\section{Purpose and audience}

The project is a shared academic resource for students taking the same university
modules. It should be readable without access to an editor's private notes and should
remain useful across different teaching years. The writing must therefore be
consistent, source-aware, and conservative in its claims.

\section{Required document structure}

Every course should contain the following three parts.

\subsection{Knowledge Notes}
Organise material according to the official syllabus and the actual teaching order.
Each chapter should normally contain:
\begin{enumerate}
  \item learning objectives;
  \item definitions and notation;
  \item main results and methods;
  \item proofs or derivations required by the module;
  \item worked examples;
  \item common errors and examination notes;
  \item a concise chapter summary.
\end{enumerate}

\subsection{TTL Materials}
Each TTL sheet should state its date or number, related topics, and source. Solutions
should explain the method rather than record only the final answer. When two methods
are genuinely useful, present the standard method first and the alternative method
second.

\subsection{Past Papers}
Each paper must preserve its academic year, assessment type, question numbering, mark
allocation, and source status. Distinguish carefully among official solutions,
editorial solutions, reconstructed questions, and unofficial marking guidance.

\section{Language and tone}

\begin{itemize}
  \item Write primarily in English, using complete sentences and standard academic
        terminology.
  \item Prefer direct statements over conversational commentary.
  \item Avoid first-person narration unless an editorial decision must be explained.
  \item Define abbreviations at first use. Retain ``TTL'' only if this is the official
        or commonly accepted course term.
  \item Do not imitate a lecturer's informal wording when a precise mathematical
        formulation is available.
\end{itemize}

\section{Mathematical writing}

\subsection{Definitions and notation}
Define every nonstandard symbol before use. Use one notation consistently throughout
the course. If the official lecture notes use conflicting notation, select one system
for the compiled notes and record the alternative in a remark.

\subsection{Theorems and proofs}
State all hypotheses explicitly. A proof should identify the main argument and include
enough intermediate steps for the intended audience. Do not add advanced generality
unless it improves the course-level explanation.

\subsection{Calculations}
Align multi-step calculations and justify nontrivial transitions. Numerical answers
should include units where appropriate. Approximate values should state the rounding
convention or number of significant figures.

\subsection{Solutions}
A complete solution should contain:
\begin{enumerate}
  \item the relevant method or theorem;
  \item the essential working;
  \item the final conclusion in a clearly identifiable form;
  \item a verification step when it is short and informative.
\end{enumerate}

\section{Source and accuracy policy}

\begin{itemize}
  \item Cite the lecture, textbook, TTL sheet, or examination paper from which material
        is taken.
  \item Never label an editorial answer as an official solution.
  \item Mark uncertain transcriptions with \CheckTag\ and resolve them before a stable
        release.
  \item Record substantive corrections in the revision history.
  \item Respect university rules governing the distribution of assessments, model
        answers, and teaching materials.
\end{itemize}

\section{Difficulty and relevance labels}

Use labels sparingly and consistently:
\begin{center}
\begin{tabular}{ll}
\CoreTag & required knowledge or a foundational method;\\[3pt]
\ExamTag & directly relevant to common assessment formats;\\[3pt]
\ProofTag & a proof that students may be expected to reproduce or adapt;\\[3pt]
\CalculationTag & a standard computational procedure;\\[3pt]
\ExtensionTag & useful enrichment beyond the minimum course requirement;\\[3pt]
\CheckTag & content requiring editorial verification.
\end{tabular}
\end{center}

\section{File and folder conventions}

Use stable, descriptive file names without spaces where possible. A recommended course
folder is:
\begin{verbatim}
course-code-course-name/
  main.tex
  course-notes.cls
  frontmatter/
    revision-history.tex
  chapters/
    knowledge-notes.tex
    ttl-materials.tex
    past-papers.tex
  figures/
  tables/
  bibliography.bib
\end{verbatim}

For a large course, split each part into one file per chapter, TTL sheet, or examination
year. Do not place generated auxiliary files in the source folders.

\section{Editorial workflow}

A stable release should pass the following sequence:
\begin{enumerate}
  \item \textbf{Draft}: one editor enters the content and sources;
  \item \textbf{Mathematical review}: another editor checks definitions, computations,
        and proofs;
  \item \textbf{Language review}: notation, terminology, and formatting are normalised;
  \item \textbf{Source review}: past-paper status and distribution permissions are
        checked;
  \item \textbf{Release}: the PDF is compiled, visually inspected, and assigned a
        version number.
\end{enumerate}

\section{Use of the template}

Edit only the metadata block in \texttt{main.tex} and the content files under
\texttt{chapters/}. Global layout changes belong in \texttt{course-notes.cls} and
should be made centrally rather than copied into individual course projects.

To create a student copy without displayed solutions, use
\begin{verbatim}
\documentclass[nosolutions]{course-notes}
\end{verbatim}
instead of the default \texttt{solutions} option.
